\usepackage[nospace]{cite}  % for smart citation
\usepackage{geometry}  % for easy margin settings
%
% margins setting
\geometry{verbose,a4paper,tmargin=3.5cm,bmargin=2cm,lmargin=4cm,rmargin=2cm}
%

% 插圖套件 graphicx
% 使用者工作流程是用 pdftex 還是 latex + dvipdfmx?
% 視情況而有不同的參數
% 這裡作自動判斷
% 參考自
% http://www.tex.ac.uk/cgi-bin/texfaq2html?label=ifpdf
%\newcommand\mydvipdfmxflow{dvipdfmx}
%\newcommand\mypdftexflow{pdftex}
%\ifx\pdfoutput\undefined
%  % not running pdftex
%  \usepackage[dvipdfm]{graphicx}
%  \newcommand\myworkflow{dvipdfmx}  % set the flag for hyperref
%\else
%  \ifx\pdfoutput\relax
%    % not running pdftex
%    \usepackage[dvipdfm]{graphicx}
%    \newcommand\myworkflow{dvipdfmx}  % set the flag
%  \else
%    % running pdftex, with...
%    \ifnum\pdfoutput>0
%      % ... PDF output
%      \usepackage[pdftex]{graphicx}
%      \newcommand\myworkflow{pdftex}  % set the flag
%    \else
%      %...DVI output
%      \usepackage[dvipdfm]{graphicx}
%      \newcommand\myworkflow{dvipdfmx}  % set the flag
%    \fi
%  \fi
%\fi

% 由於使用xelatex，所以直接搭配使用dvipdfmx
%\usepackage[dvipdfmx]{graphicx}
% 增強功能型頁楣 / 頁腳套件
\usepackage{fancyhdr}  % 借用此套件來擺放浮水印 
% (佔用了 central header)
% 不需要浮水印的使用者仍可利用此套件，產生所需的 header, footer
%
% 啟動 fancy header/footer 套件
\pagestyle{fancy}
\fancyhead{}  % reset left, central, right header to empty
\fancyfoot[C]{\thepage} %中間 footer 擺放頁碼
\renewcommand{\headrulewidth}{0pt} % header 的直線; 0pt 則無線

\setlength{\headheight}{15pt}


% 如果不需要任何浮水印，則請把下列介於 >>> 與 <<< 之間
% 的文字行關掉 (行首加上百分號)
%% 浮水印 >>> 
%%
% this file is encoded in utf-8
% v2.0 (Apr. 5, 2009)
% 如果浮水印不是全篇需要，請把下列介於 >>> 與 <<<
% 的「全篇浮水印專用碼」關掉 (行首加百分號)
% 參考自 Keith Reckdahl 寫的 "Using Imported Graphics in LATEX2e" (epslatex.pdf) p.39
% 如果只有特定頁需要浮水印
% 則依該頁屬性使用下列之一的命令 
% 普通頁命令 \thispagestyle{WaterMarkPage}
% plain 頁命令 \thispagestyle{PlainWaterMarkPage}
% empty 頁命令 \thispagestyle{EmptyWaterMarkPage}


% 將重複使用的浮水印章
% 圖檔是 my_watermark.xxx
% 副檔名可以不加，可以是 latex 系統能處裡的任何格式：pdf, gif, png, jpg, eps, ...
% 某些圖檔格式在某些工作流程可能需要作前置處裡。
% 例如，pdflatex 無法直接處理 eps 檔
%  latex + dvipdfmx 無法直接處理 pdf, gif, png, jpg, 需要用 ebb 小工具程式
%  對圖檔產生 .bb 對應檔。
%
% 寬為 5.1 cm
\newsavebox{\mywatermark}
%\sbox{\mywatermark}{\includegraphics[keepaspectratio,%
%width=5.1cm]{my_watermark}}

\sbox{\mywatermark}{\includegraphics[keepaspectratio,%
width=2.5cm]{watermark/ntust_watermark.eps}}   %watermark for NTUST by saiba

% 將 central header 擺放浮水印的巨集指令
\newcommand{\PlaceWaterMark}{\fancyhead[C]{\setlength{\unitlength}{1in}%
\begin{picture}(0,0)%
%\put(-1,-6.2){\usebox{\mywatermark}}% 圖檔擺放的位置座標 for yzu
\put(-0.5,-6.2){\usebox{\mywatermark}} % for NTUST template by saiba
\end{picture}}%
}

\fancyhead{}  % reset left, central, right header to empty
%% 如果不需整篇論文都要浮水印
%% 則下面  >>> 與 <<< 之間的程式碼請關閉
%% >>> 全篇浮水印
\PlaceWaterMark  % 每一頁都有浮水印 (除了 plain、empty 頁以外)

% 重新定義 plain 頁面
% 每張 plain 頁面 (每一章的第一頁) 也加浮水印

\fancypagestyle{plain}{%
\fancyhead{}%
\PlaceWaterMark%
\fancyfoot{}%
\fancyfoot[C]{\thepage}
\renewcommand{\headrulewidth}{0pt}%
\renewcommand{\footrulewidth}{0pt}%
}
%% <<< 全篇浮水印

%% 如果只有一、兩頁需要有浮水印
%% 可以在該頁 (有頁碼) 使用 \thispagestyle{WaterMarkPage}
%% 此命令不影響原有的 header、footer
\fancypagestyle{WaterMarkPage}{%
\PlaceWaterMark%
}

%% 如果只有一、兩頁 plain 頁需要有浮水印 (如 摘要、自傳等)
%% 可以在該頁 (有頁碼) 使用 \thispagestyle{PlainWaterMarkPage}
%% 只有頁碼與浮水印，沒有其他的 header、footer
%% 等同於 plain page style + water mark
\fancypagestyle{PlainWaterMarkPage}{%
\fancyhead{}%
\PlaceWaterMark%
\fancyfoot{}%
\fancyfoot[C]{\thepage}
\renewcommand{\headrulewidth}{0pt}%
\renewcommand{\footrulewidth}{0pt}%
}

%% 如果只有一、兩頁 empty 頁需要有浮水印 (如封面、書名頁)
%% 可以在該頁 (無頁碼) 使用 \thispagestyle{EmptyWaterMarkPage}
%% 等同於 empty page style + water mark
\fancypagestyle{EmptyWaterMarkPage}{%
\fancyhead{}%
\PlaceWaterMark%
\fancyfoot{}%
\renewcommand{\headrulewidth}{0pt}%
\renewcommand{\footrulewidth}{0pt}%
}

%% <<< 浮水印

% 如需額外的頁楣 (header) 或 footer，請在 my_headerfooter.tex 裡依例修改
% 它的預設內容是都關掉，可依需要打開
%%
% this file is encoded in utf-8
% v2.0 (Apr. 5, 2009)

%%%%%%% 其他的 header (left, right) 定義
% 底下定義了一些常見的 header 型式
% 預設情況是關掉的
% 使用者可以視需要將之打開
% 也就是把下列介於 >>> 與 <<< 之間
% 的文字行打開 (行首去掉百分號)

%% header >>>
%\renewcommand{\chaptermark}[1]{%
%\markboth{\prechaptername\ \thechapter\ \postchaptername%
%\ #1}{}%
%}  %定義 header 使用的「章」層級的戳記
%\fancyhead[L]{} % 左 header 為空
%\fancyhead[R]{\leftmark}  % 右 header 擺放「章」層級的戳記 (以 \leftmark 叫出)
%\renewcommand{\headrulewidth}{0.4pt}  % header 的直線 0.4pt; 0pt 則無線
%% <<< header

%%%%%%% 其他的 footer (left, right) 定義
% 底下定義了一些常見的 footer 型式
% 預設情況是關掉的
% 使用者可以視需要將之打開
% 也就是把下列介於 >>> 與 <<< 之間
% 的文字行打開 (行首去掉百分號)

%% footer >>>
%\fancyfoot[L]{} % 左 footer 為空
%\fancyfoot[R]{\small{YZU \LaTeX\ v2.0}} % 右 footer 擺放論文格式版本
%\renewcommand{\footrulewidth}{0.4 pt} % footer 的直線 0.4pt; 0pt 則無線
%% <<< footer




%%%%%%%%%%%%%%%%%%%%%%%%%%%%%%
%%%% 非必要的套件，但很實用
\usepackage{amsmath} % 各式 AMS 數學功能
\usepackage{amssymb} % 各式 AMS 數學符號
\usepackage{mathrsfs} %草寫體數學符號，在數學模式裡用 \mathscr{E} 得草寫 E
\usepackage{listings} % 程式列表套件
\usepackage{subfigure}
%
% listing setting
\lstset{breaklines=true,% 過長的程式行可斷行
extendedchars=false,% 中文處理不需要 extendedchars
texcl=true,% 中文註解需要有 TeX 處理過的 comment line, 所以設成 true
comment=[l]\%\%,% 以雙「百分號」做為程式中文註解的起頭標記，配合 MATLAB
basicstyle=\small,% 小號字體, 約 10 pt 大小
commentstyle=\upshape,% 預設是斜體字，會影響註解裏的英文，改用正體
%language=Octave % 會將一些 octave 指令以粗體顯示
}

\usepackage{url} % 在文稿中引用網址，可以用 \url{http://www.yzu.edu.tw} 方式

%%%% 以上為非必要套件
%%%%%%%%%%%%%%%%%%%%%%%%%%%%%%

%%% 以下是 hyperref 套件
%%%%%%%%%%%%%%%%%%%%%%%%%%%%%%
% hyperref 會擾亂 cite.sty 對文獻號碼縮編的排版，所以依據
% http://www.ctan.org/tex-archive/macros/latex/contrib/hyperref/
% 作如下的更動，使得 hyperref 不做文獻號碼的超連結。
\makeatletter
\def\NAT@parse{\typeout{This is a fake Natbib command to fool Hyperref.}}
\makeatother

% hyperlinkable table of contents
% 章節目錄、圖表超連結
%\ifx\myworkflow\mydvipdfmxflow
	%\usepackage[dvipdfmx, debug, colorlinks, linkcolor=black, citecolor=black, urlcolor=black, unicode]{hyperref} % remark by saiba 20110504 since dvipdfmx not supported in xelatex 
	\usepackage[debug, colorlinks, linkcolor=black, citecolor=black, urlcolor=black]{hyperref}  % remove dvipdfmx, unicode by saiba
%\else
%	\usepackage[pdftex, debug, colorlinks, linkcolor=black, citecolor=black, urlcolor=black]{hyperref}	
%\fi

% if hyperref is not used (e.g., in LyX application)
% define dummy \phantomsection for those occurences
%   in ntust_frontpages.tex, ntust_backpages.tex, my_appendix.tex
\ifx\hypersetup\undefined
	\newcommand\phantomsection{}
\fi
%%%% 以上為所有套件
%%%% 
%%%% 

% global page layout
\newcommand{\mybaselinestretch}{1.5}  %行距 1.5 倍 + 20%, (約為 double space)
\renewcommand{\baselinestretch}{\mybaselinestretch}  % 論文行距預設值
\parskip=2ex  % 段落之間的間隔為兩個 x 的高度
\parindent = 0Pt  % 段首內縮由 CJK 控制，所以這裡就設成不內縮

%%%%%%%%%%%%%%%%%%%%%%%%%%%%%
%  end of preamble
%%%%%%%%%%%%%%%%%%%%%%%%%%%%%